%%%%%%%%%%%%%%%%%%%%%%%%%%%%%%%%%%%%%%%%%
% Tufte Essay
% LaTeX Template
% Version 2.0 (19/1/19)
%
% This template originates from:
% http://www.LaTeXTemplates.com
%
% Authors:
% The Tufte-LaTeX Developers (https://www.ctan.org/pkg/tufte-latex)
% Vel (vel@LaTeXTemplates.com)
%
% License:
% Apache License, version 2.0
%
%%%%%%%%%%%%%%%%%%%%%%%%%%%%%%%%%%%%%%%%%

%----------------------------------------------------------------------------------------
%	PACKAGES AND OTHER DOCUMENT CONFIGURATIONS
%----------------------------------------------------------------------------------------

\documentclass[a4paper]{tufte-handout} % Use A4 paper by default, remove 'a4paper' for US letter

\usepackage{graphicx} % Required for including images
\setkeys{Gin}{width=\linewidth, totalheight=\textheight, keepaspectratio} % Default images settings
%\graphicspath{{Figures/}{./}} % Specifies where to look for included images (trailing slash required)

\usepackage{amsmath, amsfonts, amssymb, amsthm} % For math equations, theorems, symbols, etc
%\usepackage{units} % Non-stacked fractions and better unit spacing
\usepackage{physics,siunitx}
\usepackage{cancel}

\usepackage{booktabs} % Required for better horizontal rules in tables
\usepackage{subfiles}

%----------------------------------------------------------------------------------------
%	TITLE SECTION
%----------------------------------------------------------------------------------------

\title{Length Scale Analysis}

\author{Francisco Vazquez-Tavares}

\date{\today} % Date, use \date{} for no date


%----------------------------------------------------------------------------------------
%	COMMANDS SECTION
%----------------------------------------------------------------------------------------


%----------------------------------------------------------------------------------------

\begin{document}

\maketitle % Print the title section

%----------------------------------------------------------------------------------------
%	ABSTRACT/SUMMARY
%----------------------------------------------------------------------------------------

\begin{abstract}
    \justifying

    Document the length analysis of the system.
\end{abstract}

%----------------------------------------------------------------------------------------
%	ESSAY BODY
%----------------------------------------------------------------------------------------

\justifying

\section{Characteristic distances of the system}

The building blocks of the network are the interactions between patchy particles.
From this two elements, the \textit{reference} length is defined.
The central particle has a radius of \(1/2\) length units (\(\sigma_C/2\)), meanwhile the patches have a radius of \(1/5\) length units (\(\sigma_P/2\)).
The distance between the center of those particles is \(2/5\) (\(h_{\mathrm{P}}\)).
On the other hand, the distance range of the interaction (\(r_{\mathrm{b}}\)) between patches is \(\sigma_P\) length units and the cutoff distance is \(3/2\sigma_P\) length units.
With these distances we can define the distance of interaction between two patchy particles,
\begin{gather*}
    h_{\mathrm{P}} + r_{\mathrm{b}} + h_{\mathrm{P}} = 2\sigma_P + r_{\mathrm{b}}.
\end{gather*}
This distance is going to be refer as  \(r_{\mathrm{int}}\) and has the following range \(\in[6/5,7/5]\).
Then, if we increase the amount of patchy particles in pair-wise interaction, the distance between central particles is given by
\begin{gather*}
        d_{\mathrm{int}} = \left(N-1\right)\left(2\sigma_p + r_b\right).
\end{gather*}


\end{document}
